% ==============================================================
% CAPÍTULO 1 — INTRODUÇÃO
% ==============================================================

\chapter{Introdução}
\label{cap:introducao}

O presente relatório apresenta o projeto desenvolvido para a disciplina PRO3584 -- Projeto, Processo e Gestão da Inovação, ministrada pelo professor Mario Sergio Salerno no Departamento de Engenharia de Produção da Escola Politécnica da Universidade de São Paulo. O objetivo central do trabalho é aplicar os conceitos e métodos estudados na disciplina a um caso real de empresa inovadora, articulando o embasamento teórico com a prática e desenvolvendo a capacidade de diagnóstico e proposição de melhorias em sistemas de gestão da inovação.

Para esta finalidade, selecionou-se a Capim Software, healthtech--fintech fundada em 2021 com o objetivo de digitalizar a gestão financeira e operacional de clínicas odontológicas. A empresa atua estrategicamente no cruzamento entre dois ecossistemas de alto crescimento no Brasil -- healthtech e fintech -- e construiu, em pouco mais de quatro anos, uma base de mais de 6.000 clínicas atendidas e R\$~450 milhões em crédito viabilizado para tratamentos odontológicos.

\section{Objetivo e Delimitação do Escopo}

O escopo central desta pesquisa concentra-se na análise do processo decisório de inovação da Capim -- seus pontos de decisão, gates, comitês e critérios -- dentro do sistema mais amplo de gestão da inovação da organização. Uma primeira versão deste relatório privilegiou a caracterização do negócio, do mercado e da estrutura organizacional da empresa por meio de ferramentas de uso geral (Estrela de Galbraith, configuração de Mintzberg, análise SWOT). A devolutiva do professor orientador apontou que essas análises, embora úteis para contextualizar a empresa, não substituem a aplicação direta dos métodos de gestão da inovação trabalhados na disciplina ao próprio sistema decisório, e identificou uma lacuna na discussão sobre como a empresa gera e seleciona as ideias que alimentam seu funil de projetos.

Esta versão do relatório responde diretamente a essa devolutiva. A caracterização da empresa e do mercado permanece, pois contextualiza as decisões analisadas; o núcleo analítico do trabalho passa a ser, porém, a aplicação de quatro frameworks centrais da disciplina ao processo decisório da Capim: a tipologia contingencial de processos de inovação de \textcite{salerno2015}, a cadeia de valor da inovação de \textcite{hansenbirkinshaw2007}, a gestão de portfólio de risco e retorno de \textcite{goffinmitchell2010} e o modelo de ambidestria Descoberta--Incubação--Aceleração tratado por \textcite{salernogomes2018}. A partir desse diagnóstico, o relatório avalia explicitamente se a geração de ideias constitui um problema para a empresa e propõe um sistema aprimorado de gestão da inovação, incluindo um \textit{learning plan} formal para as iniciativas de maior incerteza, nos termos de \textcite{riceoconnorpierantozzi2008}.
